\section{Syntax}\label{sec:syntax}

\subsection{Grammar}

\begin{center}
\begin{tabular}{rcl}
$\varphi$ &::=& \texttt{true} $\mid$ \texttt{false} $\mid$ $p$
              $\mid$ \texttt{\{} \textit{arith} \texttt{\}} \\
          &$\mid$& \texttt{!}$\varphi$ $\mid$ $\varphi$ \texttt{\&} $\varphi$
              $\mid$ $\varphi$ \texttt{|} $\varphi$
              $\mid$ $\varphi$ \texttt{->} $\varphi$
              $\mid$ $\varphi$ \texttt{<->} $\varphi$ \\
          &$\mid$& \texttt{EF}$I\,\varphi$ $\mid$ \texttt{AF}$I\,\varphi$
              $\mid$ \texttt{EG}$I\,\varphi$ $\mid$ \texttt{AG}$I\,\varphi$ \\
          &$\mid$& \texttt{E(}$\varphi$ \texttt{U}$I$ $\varphi$\texttt{)}
              $\mid$ \texttt{A(}$\varphi$ \texttt{U}$I$ $\varphi$\texttt{)} \\
          &$\mid$& \texttt{E(}$\varphi$ \texttt{R}$I$ $\varphi$\texttt{)}
              $\mid$ \texttt{A(}$\varphi$ \texttt{R}$I$ $\varphi$\texttt{)} \\[4pt]
$I$       &::=& $\varepsilon$ $\mid$ \texttt{[}$a$\texttt{,}$b$\texttt{]}
              $\mid$ \texttt{(}$a$\texttt{,}$b$\texttt{)}
              $\mid$ \texttt{[}$a$\texttt{,}$b$\texttt{)}
              $\mid$ \texttt{(}$a$\texttt{,}$b$\texttt{]}
              $\mid$ \texttt{[}$a$\texttt{,inf)} \\
\end{tabular}
\end{center}

\noindent
Atoms are identifiers such as \texttt{p}, \texttt{q}, \texttt{p1}.  An omitted
interval means $[0,\infty)$, so \texttt{EF p} and \texttt{EF[0,inf) p} are the
same formula.

\subsection{Operators}

\begin{center}
\begin{tabular}{llll}
\toprule
Written & Means & Timed form & Reading \\
\midrule
\texttt{EF} $\varphi$ & $\exists\Diamond$ & \texttt{EF[1,5] p} & some run reaches $p$ within $[1,5]$ \\
\texttt{AF} $\varphi$ & $\forall\Diamond$ & \texttt{AF[0,3] p} & every run reaches $p$ within $[0,3]$ \\
\texttt{EG} $\varphi$ & $\exists\Box$     & \texttt{EG[1,5] p} & some run keeps $p$ throughout $[1,5]$ \\
\texttt{AG} $\varphi$ & $\forall\Box$     & \texttt{AG[0,2] p} & every run keeps $p$ throughout $[0,2]$ \\
\texttt{E(}$\varphi$ \texttt{U} $\psi$\texttt{)} & $\exists\,\mathbf{U}$ &
  \texttt{E(p U[1,3] q)} & some run holds $p$ until $q$, in $[1,3]$ \\
\texttt{A(}$\varphi$ \texttt{R} $\psi$\texttt{)} & $\forall\,\mathbf{R}$ &
  \texttt{A(p R[2,5] q)} & $\psi$ holds until released by $\varphi$ \\
\bottomrule
\end{tabular}
\end{center}

\paragraph{There is no next operator.}
TCTL has none, and \texttt{EX}/\texttt{AX} are rejected by the parser.  They
exist only as internal markers inside the construction.

\subsection{Intervals}

Bounds are non-negative integers.  Both closed and open endpoints are supported,
and the strictness is respected exactly: \texttt{EF[1,2] p} and
\texttt{EF(1,2) p} are different questions.

A \emph{punctual} interval $[a,a]$ is accepted --- see the completeness note in
Section~\ref{sec:syntax}'s preceding chapter.  An empty interval (a lower bound
above the upper one, or an open singleton) is a parse error.

\subsection{Arithmetic leaves}

Constraints in braces are handed to Z3 rather than to the temporal machinery:

\begin{lstlisting}[style=formula]
{x <= 5} & {x > 10}          % unsatisfiable: no such x
EF ({x <= 10} & p)           % satisfiable
\end{lstlisting}

\noindent
Variables are integer-sorted by default; \texttt{--real-vars x,y} makes the
named ones reals.
